\subsection{A positive repair lemma}
\label{sec:positive-repair}

The seven-qubit purity and entropy certificates
(Theorems~\ref{thm:fivecopy-sharp-certificate}
and~\ref{thm:seven-entropy-localized-identity}) both correct a strictly positive
baseline, which fails the target identity, into an exact positive identity. The
correction is found from finitely many evaluations, and its existence and
positivity follow from two verified norm bounds. The common step is the following
finite-dimensional real linear algebra; each application supplies its own spaces
and verifies the bounds.

\begin{lemma}[positive repair]
\label{lem:positive-repair}
Let $\mathcal V$ be a finite-dimensional real vector space, $\mathcal C\subseteq\mathcal V$
a subspace of dimension $d$, and $\mathcal E:\mathcal V\to\mathbb R^m$ a linear map. Let
$\Phi:\mathbb R^N\to\mathcal V$ be linear with $\Phi(\mathbb R^N)\subseteq\mathcal C$, let
$\rho\in\mathcal C$, and put $A=\mathcal E\Phi\in\mathbb R^{m\times N}$ and
$r=\mathcal E\rho\in\mathbb R^m$. Suppose $L\in\mathbb R^{d\times m}$ and
$R\in\mathbb R^{N\times d}$ satisfy
\begin{equation}\label{eq:positive-repair-inverse}
 \|I_d-LAR\|_2\le\delta<1 .
\end{equation}
\begin{enumerate}
\item[(a)] $A$ has rank $d$; $\Phi$ maps onto $\mathcal C$, and $\mathcal E$ is injective
  on $\mathcal C$.
\item[(b)] $z=R(LAR)^{-1}Lr$ satisfies $\Phi(z)=\rho$ and
  $\|z\|_2\le\|R\|_2\|Lr\|_2/(1-\delta)$. If the entries of $A,r,L,R$ lie in a
  subfield $K\subseteq\mathbb R$, then $z\in K^N$.
\item[(c)] Let $\Delta_b:\mathbb R^N\to\operatorname{Sym}_{k_b}(\mathbb R)$, $b=1,\ldots,B$, be
  linear maps with $\|\Delta_b(y)\|_2\le c\|y\|_2$ for every $y$. If
  \begin{equation}\label{eq:positive-repair-margin}
   \frac{c\,\|R\|_2\|Lr\|_2}{1-\delta}<1,
  \end{equation}
  then $I+\Delta_b(z)\succ0$ for every $b$.
\end{enumerate}
\end{lemma}
\begin{proof}
By~\eqref{eq:positive-repair-inverse} the Neumann series for $(LAR)^{-1}$ converges, so
$LAR$ is invertible with $\|(LAR)^{-1}\|_2\le1/(1-\delta)$. Hence $\operatorname{rank}A\ge d$.
Since $A=\mathcal E\Phi$ factors through the $d$-dimensional space $\mathcal C$, also
$\operatorname{rank}A\le d$. Equality forces $\Phi(\mathbb R^N)=\mathcal C$ and
$\ker\mathcal E\cap\mathcal C=0$, proving (a).

Since $AR$ has rank $d$, its image is $A(\mathbb R^N)=\mathcal E(\mathcal C)$, which
contains $r$. Write $r=ARy$; applying $L$ gives $y=(LAR)^{-1}Lr$, so $Az=r$. Thus
$\mathcal E(\Phi(z)-\rho)=0$ with $\Phi(z)-\rho\in\mathcal C$, and (a) gives $\Phi(z)=\rho$.
The norm bound follows from the bound on $\|(LAR)^{-1}\|_2$, and $z\in K^N$ because
$(LAR)^{-1}$ has entries in $K$. For (c),
$\|\Delta_b(z)\|_2\le c\|z\|_2<1$ by (b) and~\eqref{eq:positive-repair-margin}, so every
eigenvalue of $I+\Delta_b(z)$ is positive.
\end{proof}
No singular value enters the conclusion: $\delta$ and the margin bound the exact matrices,
and $L,R$ are arbitrary candidates (in both applications, binary64 singular-vector data read
as exact dyadic rationals).
